\chapter{Rotores}
\label{cap:mm-rotores}

\opening{Las palas son piezas críticas: los límites de daño se miden en milésimas de pulgada y las reparaciones no autorizadas han causado accidentes fatales.}

\section{Rotor principal: descripción}
\fuente{MM, capítulos 9 y 26}

El rotor principal tiene \textbf{dos palas totalmente metálicas} montadas al cubo mediante \textbf{bisagras de conicidad}; el cubo se monta al eje mediante una \textbf{bisagra de balanceo} (\emph{teeter}). Ambas bisagras usan rodamientos autolubricados. Los topes de caída (\emph{droop stops}) proporcionan una fricción en la bisagra de balanceo que normalmente impide que el rotor se balancee durante el arranque y la detención. Los rodamientos de cambio de paso de cada pala están en un alojamiento en la raíz, \textbf{lleno de aceite} y sellado con una funda elastomérica.

Cada pala tiene un larguero grueso de acero inoxidable en el borde de ataque, resistente a la corrosión y la erosión. Los revestimientos están pegados al larguero aproximadamente 1~in detrás del borde de ataque. \textbf{La pala debe repintarse si la pintura se erosiona hasta el metal desnudo en la línea de pegado revestimiento--larguero}: si la línea queda expuesta, la unión puede dañarse.

\begin{advertencia}
Por sus resultados potencialmente destructivos, \textbf{está prohibido usar cinta antierosión} en las palas.
\end{advertencia}

\section{Palas del rotor principal}

\subsection{Retiro}

Se necesitan \textbf{cuatro personas}: una sostiene la pala cerca de la punta y otra la raíz mientras se quita o pone el perno.

\begin{pasos}
\item Marque una pala, su posición en el cubo, su bieleta de paso y su perno y tuerca con ``X''; la opuesta con ``O''.
\item Mida y registre las holguras axiales de las bisagras de conicidad.
\item Desconecte las bieletas de los cuernos de paso; quite las chavetas y afloje las tuercas hasta apriete con los dedos.
\item Retire tuerca, arandela de empuje y suplementos del borde de fuga; levante la pala para que el colmillo del eje (\emph{tusk}) libre el tope de caída, gire el cuerno hacia abajo y saque el perno.
\end{pasos}

\begin{precaucion}
Tras quitar una pala, mantenga la otra nivelada hasta retirarla. No deje caer los muñones (dentro de los rodamientos del cubo). Los herrajes son propios de cada pala, de cada borde y de cada posición: vuelva a colocarlos en el cubo, a mano, exactamente como salieron.
\end{precaucion}

\subsection{Instalación}

\begin{pasos}
\item Identifique los números de parte del cubo y del eje de pala y los herrajes que les corresponden.
\item Verifique la fricción de balanceo (\S\,\ref{sec:friccion-bisagras}).
\item Si la holguras axiales registradas estaban fuera de tolerancia, o se cambió un rodamiento o un eje de pala, haga el cálculo de muñones y suplementos.
\item Con el cubo nivelado y los muñones en los rodamientos, inserte el eje de pala, alinee y coloque el perno desde el lado del borde de ataque; suplementos del borde de fuga (si los hay), arandela de empuje y tuerca a mano, con \textbf{antiagarrotante A257-9} solo en la rosca y la cara de la tuerca (muñones, suplementos y zona de rodamientos deben quedar limpios y secos).
\item Apriete hasta asentar, afloje hasta que ambas arandelas de empuje giren libres; ponga a cero el medidor de elongación MT122.
\item Con llaves de al menos 150~ft-lb, apriete hasta que alineen los agujeros de tuerca y perno y mida la \textbf{elongación del perno}.
\end{pasos}

\begin{center}
\begin{tabularx}{\linewidth}{@{}>{\raggedright\arraybackslash}X>{\raggedright\arraybackslash}X>{\raggedright\arraybackslash}X@{}}
\toprule
Cubo & Perno de conicidad & Elongación \\
\midrule
A154-1 & NAS630-80 (MS21250-10080) & 0,015--0,017~in \\
B370-1 & NAS632-82 (MS21250-12082) & 0,010--0,012~in \\
Bisagra de balanceo & --- & 0,015--0,017~in \\
\bottomrule
\end{tabularx}
\end{center}

Si la elongación queda fuera de rango, se instalan \textbf{perno y tuerca nuevos}, se estiran (\S\,23-33) y se taladran juntos (broca de cobalto de 0,125~in con aceite de corte; guía MT569). La chaveta nueva se instala \textbf{húmeda con imprimación} (preferentemente de cromato).

\begin{advertencia}
No estire de menos ni de más los pernos de balanceo o de conicidad: ambas cosas pueden causar su falla.
\end{advertencia}

Al terminar se hacen \emph{tracking} y balanceo (capítulo~\ref{cap:mm-reglaje}). Con palas de reemplazo, quite las tapas de punta tras la primera puesta en marcha y limpie los residuos.

\subsection{Fundas y llenado del alojamiento de paso}

La abrazadera interior A165-1 de la funda se aprieta a un \textbf{diámetro exterior de 2,330 $\pm$ 0,005~in}, sin que el hombro del labio quede atrapado bajo ella; la exterior A165-7 sobre la brida del cuerno, con el agujero del cuerno alineado con el del eje.

El alojamiento se llena con \textbf{fluido A257-4} con la herramienta MT147-1: recipiente unos 3~ft por encima del eje, llenado por el racor inferior y purga por el superior hasta que no haya burbujas (masajeando la funda, oscilando el eje y moviendo la punta). Se coloca el perno superior, se gira la pala y a los 5 minutos se buscan fugas. Pernos B289-2 con torque especial.

\section{Cubo del rotor principal}

\begin{itemize}
\item Al retirar: marque el lado de la tuerca del perno de balanceo y el lado del brazo de cuerda del eje; mida y registre la fricción de balanceo; vuelva a armar los herrajes en el cubo exactamente como salieron.
\item El \textbf{balanceo de cuerda} del rotor se ajusta con los muñones A106 y suplementos A117: alterar la pila desbalancea.
\item Al instalar: limpie y seque; reemplace muñones o arandelas con cromo saltado, corrosión o surcos que atraviesen el cromo. Si la fricción registrada era $<$5 o $>$15~ft-lb, si se cambió un rodamiento o no hay datos, haga el cálculo de suplementos. Cabeza del perno del mismo lado que el brazo de cuerda. Los muñones no deben ``pellizcar'' los topes de caída.
\item Rodamientos del cubo: se extraen e instalan con prensa y las herramientas MT329 (siempre con soporte al instalar los de conicidad), con imprimación húmeda (sin tocar el forro de teflón), costura del forro hacia arriba y un filete de imprimación en ambos bordes.
\end{itemize}

\subsection{Cálculo de muñones y suplementos}

\begin{center}
\small
\begin{tabularx}{\linewidth}{@{}>{\raggedright\arraybackslash}X>{\raggedright\arraybackslash}X>{\raggedright\arraybackslash}X@{}}
\toprule
Parte & Medida & Uso \\
\midrule
A106-1 / -2 / -3 / -6 & 1,000 / 0,995 / 0,990 / 1,005~in & Conicidad, cubo A154-1 (sin suplementos) \\
A106-4 / -5 & 1,300 / 1,260~in & Balanceo, con suplementos \\
A106-7 & 1,005~in & Conicidad, cubo B370-1 \\
A117-5 / -6 / -7 / -8 & 0,012 / 0,015 / 0,020 / 0,025~in & Balanceo \\
A117-48 / -49 / -50 / -51 & 0,012 / 0,015 / 0,020 / 0,025~in & Conicidad, cubo B370-1 \\
\bottomrule
\end{tabularx}
\end{center}

\textbf{Bisagra de balanceo}: espacio vacío = ancho del cubo entre caras de rodamientos $-$ ancho del eje en el perno. Bajo la cabeza: A117-7 + A106-4, o A117-6/-7/-8 + A106-5. Reste esos espesores y el muñón del lado de la tuerca; \textbf{sume 0,005--0,008~in} por la compresión de apriete: el resultado es la pila inicial de suplementos entre muñón y arandela del lado de la tuerca. Use tantos tamaños distintos como sea posible para facilitar el desplazamiento del cubo al balancear.

\textbf{Bisagra de conicidad}: espacio vacío = ancho entre caras de rodamientos $-$ ancho del eje de pala; en el B370-1 reste dos muñones A106-7 y \textbf{sume 0,003--0,006~in} para obtener la pila de suplementos; en el A154-1, sume 0,003--0,006~in y elija muñones A106-1/-2/-3/-6 cuya suma iguale el resultado (el mismo número bajo ambas cabezas).

\begin{precaucion}
Una pila inicial menor que la calculada puede dañar las arandelas de empuje y los rodamientos del cubo durante la instalación.
\end{precaucion}

\subsection{Ajuste de fricción de las bisagras}
\label{sec:friccion-bisagras}

\begin{itemize}
\item \textbf{Balanceo}: quitar suplementos del lado de la tuerca \emph{aumenta} la fricción; agregarlos la \emph{reduce}. Se mide con la herramienta MT354 en los rodamientos de conicidad y una balanza (fuerza en movimiento, no de arranque). \textbf{Fricción de balanceo: menos de 15~ft-lb.} (Un valor registrado fuera de 5--15~ft-lb al desmontar obliga a rehacer el cálculo de suplementos.)
\item \textbf{Conicidad}: en el B370-1 se varía la pila de suplementos del borde de fuga; en el A154-1, la longitud del muñón del borde de fuga (más corto, más fricción). La fricción es cero cuando hay holgura axial medible. \textbf{Holgura total por bisagra: 0,002--0,006~in} (galga entre arandela y cara del rodamiento, en cabeza y tuerca).
\item Prueba funcional: con las palas levantadas de los topes, sosteniendo una nivelada y conificando la otra, \textbf{el cubo no debe balancearse}.
\end{itemize}

\textbf{Desplazamiento del cubo} (balanceo de cuerda): con dos personas conificando las palas, se empuja el perno de balanceo con otro perno y se pasan suplementos de un lado al otro según la carta de balanceo.

\section{Inspección de las palas del rotor principal}

\begin{nota}
Las palas son piezas críticas según 14 CFR 27.602. Notifique a RHC los vacíos que excedan los límites (número de serie de pala y helicóptero, tiempo en servicio, ubicación y tamaño). Estos criterios aplican a daños posteriores a la fabricación; si sospecha una irregularidad de fabricación, consulte a RHC.
\end{nota}

\begin{precaucion}
Una pala puede repararse más de una vez, pero nunca puede superarse en un mismo lugar el máximo de material removido ni de profundidad de abolladura. \textbf{Toda pala en cuya unión pegada haya entrado líquido penetrante es chatarra.}
\end{precaucion}

El daño se mide con regla y galga de espesores, con la regla paralela a los bordes; con la pala instalada, use la regla más corta posible (una larga da lecturas falsas por la flexión natural). El material removido se estima comparando con calibre o micrómetro antes y después.

\subsection{Rayas y corrosión}

\begin{center}
\small
\begin{tabularx}{\linewidth}{@{}>{\raggedright\arraybackslash}X>{\raggedright\arraybackslash}X>{\raggedright\arraybackslash}X@{}}
\toprule
Daño (después del retrabajo) & A016-6 & A016-4 \\
\midrule
Rayas a más de 15° del eje de envergadura & 0,004~in & 0,002~in \\
Rayas a menos de 15° & 0,006~in & 0,003~in \\
Corrosión RS~127,0--151,1 & 0,012~in & --- \\
Corrosión RS~75,5--127,0 & 0,008~in & --- \\
Corrosión RS~21,0--75,5 & 0,006~in & --- \\
Corrosión RS~10,4--21,0 & 0,010~in & --- \\
\bottomrule
\end{tabularx}
\end{center}

Rayas y corrosión de más de 0,0005~in se pulen con radio de transición de 0,10~in.

\subsection{Abolladuras}

Golpee en la zona (prueba de golpeteo) con una arandela AN970-4 o una moneda de 25~centavos de EE.\,UU.\ de 1965 o posterior; si hay vacíos asociados, consulte a RHC. No repare abolladuras con corte o rotura del revestimiento; la abolladura debe tener \textbf{radio de fondo mínimo de 0,060~in}.

\begin{center}
\small
\begin{tabularx}{\linewidth}{@{}>{\raggedright\arraybackslash}p{0.38\linewidth}>{\raggedright\arraybackslash}p{0.24\linewidth}>{\raggedright\arraybackslash}X@{}}
\toprule
Zona & A016-6 & A016-4 \\
\midrule
Panal: abultamiento opuesto & 0,020~in & 0,020~in \\
Panal: abolladura (exterior) & 0,125~in (RS~127,0--151,1) & 0,125~in (RS~130,3--151,1) \\
Panal: abolladura (media) & 0,090~in (RS~75,5--127,0) & 0,090~in (RS~64,0--130,3) \\
Panal: abolladura (interior) & 0,030~in (RS~10,4--75,5) & 0,030~in (RS~13,7--64,0) \\
Borde de ataque de dobladores & 0,010~in & --- \\
Uniones pegadas apoyadas & 0,006~in & 0,002~in (solo tapa de punta) \\
Deformación local (0,75~in del borde de fuga) & 0,060~in (RS~30--151); 0,015~in (RS~10,4--30) & 0,060~in (RS~30--151); 0,015~in (RS~13,7--30) \\
Larguero (radio de transición $\ge$1,0~in) & 0,020~in (RS~75,5--151); 0,010~in (RS~21--75,5) & 0,020~in (RS~64--151); 0,010~in (RS~21--64) \\
\bottomrule
\end{tabularx}
\end{center}

En el larguero no se permite pulir a menos de 0,010~in del borde de ataque de la ranura. Las abolladuras lisas de fondo redondeado pueden rellenarse y carenarse.

\widefigure[0.9\linewidth]{assets/mm/pala-abolladuras.png}{Límites de abolladuras y deformaciones locales en las palas A016-6 (izquierda) y A016-4 (derecha). Las tramas indican la profundidad máxima admisible en cada zona; RS es la estación radial en pulgadas. Fuente: MM, figura 26-12.\label{fig:pala-abolladuras}}

\subsection{Raíz de la pala}

\begin{center}
\small
\begin{tabularx}{\linewidth}{@{}>{\raggedright\arraybackslash}X>{\raggedright\arraybackslash}X>{\raggedright\arraybackslash}X@{}}
\toprule
Zona & A016-6 & A016-4 \\
\midrule
Radio de transición mínimo & 0,030~in & 0,10~in \\
Cara interior de la brida & 0,002~in & 0,010~in \\
Cara exterior de la brida & 0,005~in, diám.\ 0,250~in, máx.\ 3, a $\ge$0,10~in de agujeros & 0,005~in en 0,250~in fuera de la brida \\
Zonas expuestas del herraje & 0,040~in & 0,060~in \\
\bottomrule
\end{tabularx}
\end{center}

\widefigure[0.85\linewidth]{assets/mm/pala-raiz.png}{Límites de pulido en el herraje de raíz. Fuente: MM, figura 26-13.\label{fig:pala-raiz}}

\subsection{Vacíos (despegues)}

\begin{advertencia}
Los vacíos o despegues en las palas \textbf{no son reparables en campo}.
\end{advertencia}

\begin{itemize}
\item \textbf{Áreas críticas}: a menos de 0,50~in en envergadura y 0,30~in en cuerda del borde de cualquier unión estructural (y toda área no definida como semicrítica o no crítica). Vacío máximo \textbf{0,10~in²}; al menos 90\,\% pegado. Vacíos separados menos de 0,25~in se consideran continuos.
\item \textbf{Semicríticas} (solo A016-4, alrededor de la aleta compensadora): vacío máx.\ círculo de 0,80~in o 0,90~in²; máx.\ 0,10~in² dentro de un área crítica; 80\,\% pegado.
\item \textbf{No críticas}, A016-6: dobladores --- 80\,\% pegado; máx.\ 5,0~in², 1,0~in en cuerda y 6,0~in en envergadura; se permiten vacíos a 0,30~in del borde de ataque del doblador donde envuelve larguero y raíz. Panal --- 80\,\% pegado; máx.\ 8,0~in² (1,0 $\times$ 20,0~in) por dentro de RS~92 y 14,0~in² (2,0 $\times$ 20,0~in) por fuera.
\item \textbf{No críticas}, A016-4: dobladores A934 --- máx.\ 5,0~in² (1,10 $\times$ 6,0~in), 2,0~in desde las puntas (SL-55); panal --- 8,0~in² por dentro de RS~106 y 14,0~in² por fuera; entre RS~120--133,5 y 150--166, a $\ge$1,0~in del borde de fuga del panal y sin que el revestimiento se mueva al flexionar las aletas.
\end{itemize}

\widefigure[0.85\linewidth]{assets/mm/pala-adherencia.png}{Áreas de unión crítica, semicrítica y no crítica. Fuente: MM, figura 26-14.\label{fig:pala-adherencia}}

\section{Reparación de las palas del rotor principal}

\begin{advertencia}
Las reparaciones no autorizadas de palas han causado accidentes fatales. \textbf{No use herramientas eléctricas, decapantes químicos ni removedores químicos de corrosión.}
\end{advertencia}

\begin{pasos}
\item Mida el daño.
\item Elimine daños en bordes de fuga, aletas, tapa de punta o esquina de punta recortando dentro de los límites.
\item Pula con lija de óxido de aluminio al agua grano 220 o más fino y termine con 320 o más fino (se permite lima fina a lo largo del larguero y del borde de fuga, terminando con 320). Lije en sentido de la envergadura, quite solo lo necesario y mantenga el borde de fuga escuadrado.
\item Mida el material removido y verifique los límites.
\item Selle o rellene: limpie con solvente aprobado; imprimación (de cromato preferida) en uniones con poros, \textbf{24~h} mínimo de curado; lije el adhesivo curado a forma aerodinámica sin quitar metal; lije la pintura circundante hasta que quede un 25\,\% de imprimación.
\item Pinte y haga \emph{tracking} y balanceo según corresponda.
\end{pasos}

\textbf{Recorte}: se permite en el borde de fuga de los revestimientos y bordes de las aletas, y en tapa y esquina de punta, dentro de los límites de la figura (como alternativa, una mella del borde de fuga puede difuminarse 1,0~in mínimo a cada lado). \textbf{No se permite recortar larguero ni dobladores.} Cuerda mínima: 7,60~in en la punta de la A016-6 y 7,10~in en el resto. Lime el borde de fuga escuadrado, nunca en punta.

\widefigure[0.85\linewidth]{assets/mm/pala-recorte.png}{Límites de recorte de las palas. Fuente: MM, figura 26-15.\label{fig:pala-recorte}}

\subsection{Pintura y borde de ataque}

\begin{itemize}
\item Al secar la pintura con calor, \textbf{no supere 175~°F} en la superficie de la pala.
\item Imprimación: dos manos (cromato preferido), escarificando la primera con lija 320 o Scotch-Brite muy fino. Esmalte de poliuretano gris oscuro (raíz), negro mate, blanco y amarillo según el esquema. Deje curar la pintura Imron \textbf{al menos 72~h} antes de volar con llovizna, lluvia o polvo.
\item \textbf{Borde de ataque sin pintar}: se quita la oxidación con Scotch-Brite granate 7447, se limpia con jabón suave y agua (acetona si hace falta; \emph{no} MEK ni alcohol) y se aplica ProtectaClear con guante de nitrilo (seca en una hora; puede darse una segunda mano).
\item Para decapar el borde de ataque exterior se traza una línea a 0,65~in del borde, desde la punta hasta 72~in hacia adentro, y se lija a mano con 220 sin quitar metal; promotor de adherencia (PreKote), imprimación, pintura negra, y ProtectaClear sobre el metal desnudo.
\end{itemize}

\subsection{Mantenimiento de las puntas}

Tras quitar las tapas redondeadas, inspeccione con lupa de 10$\times$: pintura suelta o ampollada, polvo blanco de corrosión o picaduras detrás de las líneas de pegado. Si hay metal desnudo o corrosión: lije a mano la tapa y los bordes de los revestimientos (taco plano y duro, 220 y luego 320, en sentido de la cuerda en la cara vertical y de la envergadura en las superficies con radio de 0,1~in), limpie, selle las uniones expuestas con \textbf{sellador B270-1} (polisulfuro), dos manos de imprimación epoxi y pintura amarilla entre 2 y 48~h después (curar 48~h antes de volar). Tornillos de las tapas con antiagarrotante A257-9 y torque especial, tapas al ras del perfil.

\section{Rotor de cola}
\fuente{MM, capítulo 28}

El rotor de cola tiene \textbf{dos palas metálicas} y un \textbf{cubo basculante con ángulo de conicidad fijo}. Los rodamientos de paso en la raíz tienen camisa autolubricada; los de la bisagra de balanceo son \textbf{elastoméricos}. Las palas tienen revestimientos de aluminio y herrajes de raíz forjados; mantener la pintura reduce corrosión y erosión.

\widefigure[0.8\linewidth]{assets/mm/rc-conjunto.png}{Instalación del conjunto del rotor de cola: cubo A062-1, palas A029-2, cojinetes elastoméricos, separador largo de preconicidad y pesos de balanceo. Fuente: MM, figura 28-1.\label{fig:rc-conjunto}}

\subsection{Retiro e instalación}

\begin{itemize}
\item Rotule cada bieleta con el número de serie de su pala y conserve sus herrajes (los pernos pala--bieleta pueden diferir en largo y arandelas por el balanceo dinámico). Marque la orientación del perno de balanceo.
\item Haga el balanceo estático si no se conocen los herrajes de balanceo, si se cambiaron palas o si un retrabajo cambió la masa.
\item Limpie el eje de salida y las superficies de los separadores con jabón suave (pH 7--9) y agua tibia, enjuague y seque.
\item El rotor gira \textbf{en sentido antihorario visto desde la izquierda} y las palas se conifican \textbf{hacia la caja de cola}.
\item Apriete el perno de balanceo hasta que los separadores metálicos de los cojinetes toquen el eje, sin dar torque. \textbf{Tracking}: con una guía de chapa de aluminio de 1$\times$12~in doblada a 90° a 2~in de un extremo pegada al cono de cola, marque el paso del agujero de drenaje de cada punta y bascule el rotor hasta que ambas pasen por el mismo punto \textbf{dentro de 0,125~in}; luego torque especial y verificación.
\item Al bascular, perno, separadores, arandelas y tuercas deben quedar inmóviles. Tope de uretano A119-1 con arandela A141-14 y tuerca con torque especial. Balanceo dinámico (capítulo~\ref{cap:mm-reglaje}).
\end{itemize}

\subsection{Balanceo estático}

Se hace en aire calmo con la barra MT179-4 (perpendicular al cubo) en un soporte Anderson 20 o equivalente; en los cuernos de paso se simula la bieleta con perno NAS6604, arandela A214-3 y rodamiento B115-1.

\begin{itemize}
\item \textbf{En cuerda} (rotor vertical): si la pala de arriba cae con el borde de ataque primero, es la pesada; si cae con el de fuga, es la liviana. Se corrige con la longitud del perno NAS6604-16 a -21 y arandelas del lado de la tuerca en el cuerno de la pala liviana (respetando el enganche de rosca). Tolerancia: \textbf{una arandela delgada de acero}.
\item \textbf{En envergadura} (rotor horizontal): la que cae es la pesada. Se corrige con la masa de las arandelas bajo la tuerca del perno exterior pala--cubo de la pala liviana: \textbf{cuatro arandelas}, las más grandes junto al cubo. Tolerancia: \textbf{una arandela delgada de aluminio}.
\end{itemize}

\subsection{Palas}

\begin{precaucion}
Las palas A029 son un \textbf{par emparejado} de fábrica: si se cambia solo una, pida a RHC una pala compatible con el número de serie de la que queda. El separador A137-2 crea la preconicidad y va del lado exterior de la pala, en el perno exterior.
\end{precaucion}

Los bujes A138-1 se prensan con imprimación húmeda al ras del lado interior del brazo interior del cubo (se asientan al dar torque). Torque estándar y pintura de torque solo en las tuercas interiores.

\widefigure[0.8\linewidth]{assets/mm/rc-pala-criterios.png}{Criterios de inspección de las palas del rotor de cola: abolladuras (izquierda) y rayas o corrosión (derecha). Fuente: MM, figura 28-4.\label{fig:rc-pala}}

\subsection{Inspección y reparación de las palas}

\begin{advertencia}
Toda pala del rotor de cola dañada que no pueda repararse dentro de estos límites debe retirarse de servicio de inmediato y marcarse ``chatarra''. No use herramientas eléctricas ni productos químicos.
\end{advertencia}

\begin{center}
\small
\begin{tabularx}{\linewidth}{@{}>{\raggedright\arraybackslash}p{0.5\linewidth}>{\raggedright\arraybackslash}X@{}}
\toprule
Daño & Límite \\
\midrule
Rayas o corrosión por dentro de RS~13,1 & 0,004~in \\
Rayas o corrosión por fuera de RS~13,1 & 0,008~in \\
Abolladura a menos de 0,75~in del borde de ataque & 0,010~in (solo relleno) \\
Abolladura entre 0,75~in del b.\,a.\ y 0,40~in del b.\,f. & 0,030~in (solo relleno) \\
Abolladura a menos de 0,40~in del borde de fuga & 0,030~in (enderezar a 0,015~in y rellenar) \\
Abolladura sobre la unión revestimiento--raíz & 0,006~in \\
Tapa de punta y 0,775~in desde la punta & sin abolladuras \\
Abultamiento opuesto a una abolladura & máx.\ 0,010~in \\
Mella del borde de fuga & 0,050~in; cuerda mín.\ 3,950~in; transición $\ge$1,0~in a cada lado con radio $\ge$2,0~in \\
Raíz: radio de 0,5~in alrededor de los rodamientos & sin reparaciones \\
Raíz: superficies de apriete del cuerno & 0,002~in por lado (una sola vez); paralelas y perpendiculares al agujero dentro de 0,002~in \\
Raíz: agujero de 0,250~in & máx.\ 0,252~in \\
Raíz: otras zonas expuestas & 0,002~in \\
\bottomrule
\end{tabularx}
\end{center}

\begin{itemize}
\item Las abolladuras reparables deben ser lisas, de fondo redondeado con radio $\ge$0,060~in; inspeccione con lupa 10$\times$ y retire la pala si hay metal fisurado.
\item Reemplace toda pala cuya erosión haya deformado u ondulado el borde de ataque.
\item Rayas y corrosión se difuminan en envergadura con radio $\ge$0,10~in (220 y luego 320).
\item Repintado: imprimación de conversión y epoxi; Corlar 13580S para reconstruir el perfil; dos manos de Desoprime CA7502 (entre 10~min y 8~h entre manos); esmalte Imron blanco y franjas negras.
\item \textbf{Línea de pegado de la tapa de punta}: la pintura ampollada puede indicar corrosión. Quite la pintura suelta con la uña o una espátula plástica y examine con lupa 10$\times$: se admite contacto metal--metal, pero \textbf{cualquier hueco por falta de adhesivo azul o marrón, o cualquier corrosión en la línea, exige reemplazar la pala}. Un agujero que atraviese el revestimiento también.
\item Balancee el rotor de cola tras eliminar corrosión o pintar. En ambiente corrosivo, lávelo a diario; si no puede encerar, pase WD-40 estándar por los bordes de ataque (\emph{no} ACF-50, ni WD-40 ``Specialist'', ni Salt-Away).
\end{itemize}

\subsection{Cubo del rotor de cola}

Los cojinetes elastoméricos tienen un leve ajuste a presión en el cubo y se cambian con las herramientas MT556-17 o MT556-1, con imprimación húmeda en el fondo del alojamiento y una línea en el diámetro exterior. \textbf{Si hay \emph{fretting} en los alojamientos, el cubo no es aeronavegable.} En la inspección, limpie y busque daños, desgaste y corrosión; verifique brazos rectos, agujeros sin ovalizar y alojamientos sin \emph{fretting} ni excoriación. No se admite corrosión en superficies de apriete; en zonas sospechosas, quite la pintura localmente y haga líquidos penetrantes.
